%! Logarithmic Functions — Unit 5, Lesson 5.3 — Homework
\documentclass[10pt]{article}
\usepackage{precalculus-article}
\usepackage{precalculus-boxes}
\newcommand{\TallMath}[1]{$\displaystyle #1\rule[-1.4em]{0pt}{3.2em}$}

\begin{document}

\pageheader{Unit 5, Lesson 5.3}{Homework}
\namedateperiod

\vspace{0.1in}

\begin{enumerate}[leftmargin=1.8em, itemsep=14pt]

\item For each function, state the domain, range, and vertical asymptote.

\begin{enumerate}[leftmargin=1.5em, itemsep=10pt, label=\alph*.]
  \item $f(x) = \log_5 x$

  \vspace{4pt}
  Domain: \blank{3.5cm} \quad Range: \blank{3.5cm} \quad Asymptote: $x = $ \blank{1.5cm}

  \item $g(x) = -2\log_3 x$

  \vspace{4pt}
  Domain: \blank{3.5cm} \quad Range: \blank{3.5cm} \quad Asymptote: $x = $ \blank{1.5cm}

  \item $h(x) = \tfrac{1}{2}\log_{10} x$

  \vspace{4pt}
  Domain: \blank{3.5cm} \quad Range: \blank{3.5cm} \quad Asymptote: $x = $ \blank{1.5cm}
\end{enumerate}

\item For $g(x) = -\log_2 x$:

\begin{enumerate}[leftmargin=1.5em, itemsep=8pt, label=\alph*.]
  \item Fill in the end behavior.

  \vspace{4pt}
  $\displaystyle\lim_{x \to 0^+} (-\log_2 x) = $ \blank{3cm} \qquad
  $\displaystyle\lim_{x \to \infty} (-\log_2 x) = $ \blank{3cm}

  \item Is $g$ increasing or decreasing? Use the inverse relationship to explain.

  \writelines{2}
\end{enumerate}

\item A function $h$ has the table below, and $h(x) = f(x + 2)$ for some function $f$.
      The output values increase by 1 each row. Determine whether $f$ is logarithmic.

\vspace{6pt}
\begin{center}
\renewcommand{\arraystretch}{1.8}
\begin{tabular}{|c|c|c|}
\hline
\rowcolor{sky}
$x$ (input to $h$) & $h(x)$ & $x + 2$ (input to $f$) \\
\hline
$-1$ & $0$ & $1$ \\
\hline
$1$  & $1$ & $3$ \\
\hline
$7$  & $2$ & $9$ \\
\hline
$25$ & $3$ & $27$ \\
\hline
$79$ & $4$ & $81$ \\
\hline
\end{tabular}
\end{center}

\vspace{6pt}
Are the inputs of $f$ (column 3) proportional over equal output intervals?

\vspace{4pt}
Ratios: \blank{6cm}

\vspace{4pt}
Is $f$ logarithmic? \blank{2cm} \quad What base? $b = $ \blank{2cm}

\item \textbf{Conceptual question.} Explain in your own words why $f(x) = a\log_b x$ has
      no local maximum or minimum on its natural domain. Reference EK 2.11.A.2.

\writelines{3}

\item \textbf{AP-style multiple choice.} Which of the following is \textbf{NOT} a
      characteristic of $f(x) = \log_3 x$?

\vspace{6pt}
\begin{enumerate}[leftmargin=2em, itemsep=4pt, label=(\Alph*)]
  \item The domain is $x > 0$.
  \item The range is all real numbers.
  \item The function has a local maximum at $x = 1$.
  \item The graph has a vertical asymptote at $x = 0$.
\end{enumerate}

\end{enumerate}

\vspace{0.12in}

\begin{extensionbox}
\textbf{Extension --- optional.}

Compare the end behavior of $f(x) = \log_2 x$ (base $b = 2 > 1$) and $h(x) = \log_{1/2} x$
(base $b = \tfrac{1}{2}$, where $0 < b < 1$).

\vspace{6pt}
For each function, fill in: \quad $\displaystyle\lim_{x \to 0^+} = $ \blank{2.5cm}, \quad
$\displaystyle\lim_{x \to \infty} = $ \blank{2.5cm}

\vspace{6pt}
$f$: \quad $\displaystyle\lim_{x \to 0^+}\log_2 x = $ \blank{2cm} \quad
$\displaystyle\lim_{x \to \infty}\log_2 x = $ \blank{2cm}

\vspace{6pt}
$h$: \quad $\displaystyle\lim_{x \to 0^+}\log_{1/2} x = $ \blank{2cm} \quad
$\displaystyle\lim_{x \to \infty}\log_{1/2} x = $ \blank{2cm}

\vspace{6pt}
Explain how the value of the base ($b > 1$ vs.\ $0 < b < 1$) determines whether the function
is increasing or decreasing, and how it affects the direction of the end behavior.

\writelines{3}
\end{extensionbox}

\vspace{0.1in}

\begin{reflectionbox}
What is one key characteristic of logarithmic functions that surprised you, and one question you
still have about domain, end behavior, or the identification test?

\writelines{2}
\end{reflectionbox}

\end{document}
